\documentclass{article}
\usepackage[T1]{fontenc}
\usepackage{lmodern}
\usepackage{graphicx}
\usepackage{amsmath,amssymb}
\usepackage[a4paper,margin=2cm]{geometry}
\usepackage[absolute,overlay]{textpos}
\setlength{\TPHorizModule}{1mm}
\setlength{\TPVertModule}{1mm}
\usepackage[indonesian]{babel}
\usepackage{hyperref}
\setlength{\emergencystretch}{2em}
% unit: Halaman 2

\begin{document}

% === Halaman 2 (python/native) ===

2

Pendahuluan

•

Kriptografi modern adalah era kriptografi setelah penemuan komputer digital (awal 

tahun 1970-an).

•

Perkembangan teknologi komputer digital membuat hampir semua bidang ilmu 

pengetahuan berkebang pesat, termasuk kriptografi.

•

Komputer digital merepresentasikan data dan informasi dalam biner.

•

Algoritma kriptografi modern beroperasi dalam mode bit (bandingkan dengan algoritma 

kriptografi klasik beroperasi dalam mode karakter)


→ kunci, plainteks, cipherteks,  diproses dalam rangkaian bit


→ operasi xor paling banyak digunakan di dalam algoritmanya

\end{document}
